\subsubsection{Intertwining and the General Similarity Transform}
Compatibility between two descriptions of the same algebraic action is
expressed by the starting relation \eqref{eq:transform-intertwining}. A
primed carrier is carried into the unprimed description by \(\pv{T}\). If an algebraic
action is represented by \(\pv{Z}'\) before that identification and by \(\pv{Z}\)
afterward, then commutation of the action with the change of description is required:
\begin{equation*}
\pv{T}(\pv{Z}'\pv{\psi}')=\pv{Z}(\pv{T}\pv{\psi}')
\qquad\text{for every carrier }\pv{\psi}'.
\end{equation*}
The relation \(\pv{T}\pv{Z}'=\pv{Z}\pv{T}\) is obtained by removing the
arbitrary carrier. Because \(\pv{T}\) is invertible, its unique solution is
\begin{equation*}
\pv{Z}'=\pv{T}^{-1}\pv{Z}\pv{T}.
\end{equation*}
Thus Eq.\ \eqref{eq:general-transform} is not an additional dynamical law:
it is the change-of-description law forced by the intertwining requirement.
Under the opposite carrier direction, \(\pv{T}\) is replaced everywhere by
\(\pv{T}^{-1}\), and the equivalent convention
\(\pv{Z}'=\pv{T}\pv{Z}\pv{T}^{-1}\) is obtained. Under the inverse-first
convention, the displayed positive sign in Rodrigues's formula is produced
by the positive exponent in Eq.\ \eqref{eq:rotor-element}.

The fully general complex paravectors are written as
\begin{align*}
\pv{Z}&=S+\vct{V}\cdot\sigv,
&S&\in\C,
&\vct{V}&\in\C^3,\\
\pv{T}&=f+\vct{g}\cdot\sigv,
&f&\in\C,
&\vct{g}&\in\C^3,
\end{align*}
where the dot and cross products are extended complex-bilinearly. The central
determinant of \(\pv{T}\) is \(\detf{\pv{T}}=f^2-\vct{g}^2\). When
\(\detf{\pv{T}}\neq0\), the inverse is given by the general formula
\eqref{eq:paravector-inverse}:
\begin{equation*}
\pv{T}^{-1}=\frac{\cherm{\pv{T}}}{\detf{\pv{T}}}
=\frac{f-\vct{g}\cdot\sigv}{f^2-\vct{g}^2}.
\end{equation*}
Because the scalar part \(S\) is central,
\begin{align*}
\pv{Z}'
&=S+\frac{(f-\vct{g}\cdot\sigv)
 (\vct{V}\cdot\sigv)(f+\vct{g}\cdot\sigv)}{\detf{\pv{T}}}.
\end{align*}
The first multiplication is
\begin{equation*}
(f-\vct{g}\cdot\sigv)(\vct{V}\cdot\sigv)
=-\vct{g}\cdot\vct{V}
+\bigl(f\vct{V}-i\,\vct{g}\times\vct{V}\bigr)\cdot\sigv.
\end{equation*}
After multiplication by the final factor, the scalar terms are canceled and the
remaining vector terms are
\begin{align}
\pv{Z}'&=S+\vct{V}'\cdot\sigv,\notag\\
\vct{V}'
&=\frac{(f^2+\vct{g}^2)\vct{V}
-2(\vct{g}\cdot\vct{V})\vct{g}
-2if(\vct{g}\times\vct{V})}{f^2-\vct{g}^2}.
\label{eq:similarity-components}
\end{align}
This is the similarity law for every complex paravector \(\pv{Z}\) and every
invertible complex paravector \(\pv{T}\). The complex scalar part is fixed.
Neither Eq.\ \eqref{eq:general-transform} nor
Eq.\ \eqref{eq:similarity-components} is changed when \(\pv{T}\) is multiplied
by a nonzero central complex scalar.

Products, inverses, and functions are governed by the same intertwining. For a
second algebra element \(\pv{W}\),
\begin{align}
(\pv{Z}\pv{W})'&=\pv{T}^{-1}\pv{Z}\pv{W}\pv{T}
=\pv{T}^{-1}\pv{Z}\pv{T}\,\pv{T}^{-1}\pv{W}\pv{T}=\pv{Z}'\pv{W}',\notag\\
(\pv{Z}')^{-1}&=\pv{T}^{-1}\pv{Z}^{-1}\pv{T}
\qquad\text{when }\pv{Z}^{-1}\text{ exists},\notag\\
(\pv{Z}')^n&=\pv{T}^{-1}\pv{Z}^n\pv{T},\notag\\
h(\pv{Z}')&=\pv{T}^{-1}h(\pv{Z})\pv{T}.
\label{eq:similarity-covariance}
\end{align}
The last line is valid for every convergent power-series function \(h\), and
therefore in particular
\begin{equation*}
\exp(\pv{Z}')=\pv{T}^{-1}\exp(\pv{Z})\pv{T}.
\end{equation*}
The same relation is obeyed by logarithms, roots, and rational functions
wherever the same branch and domain are defined.

Commutation of the adjugate with this similarity is also obtained. Since
\(\cherm{\pv{T}}=\detf{\pv{T}}\pv{T}^{-1}\) and
\(\cherm{(\pv{T}^{-1})}=\pv{T}/\detf{\pv{T}}\),
\begin{align*}
\cherm{\pv{Z}'}
&=\cherm{\pv{T}}\,\cherm{\pv{Z}}\,\cherm{(\pv{T}^{-1})}
=\pv{T}^{-1}\cherm{\pv{Z}}\pv{T},\\
\pv{Z}'\cherm{\pv{Z}'}
&=\pv{T}^{-1}(\pv{Z}\cherm{\pv{Z}})\pv{T}=\pv{Z}\cherm{\pv{Z}}.
\end{align*}
Thus the determinant form \(S^2-\vct{V}^2\), the spectrum, and the inverse
framework are preserved by similarity for every invertible \(\pv{T}\). The
positive coefficient norm is a different construction and is not generally
preserved by a nonunitary similarity.

Algebra elements whose action is represented by multiplication are transported
by the similarity law; this law is not automatically the Lorentz law for every
physical representation. A real event paravector is characterized by
\(\pv{X}^{\mathsf H}=\pv{X}\). Such paravectors are
spanned by Hermitian carrier bilinears
\(\pv{\psi}\pv{\psi}^{\mathsf H}\): concretely, in the Pauli realization every
Hermitian two-by-two matrix is a real linear
combination of bilinears of this form. The carrier identification is
\(\pv{\psi}'=\pv{T}^{-1}\pv{\psi}\), and therefore
\begin{equation*}
\pv{\psi}'\pv{\psi}'^{\mathsf H}
=\pv{T}^{-1}(\pv{\psi}\pv{\psi}^{\mathsf H})(\pv{T}^{-1})^{\mathsf H}.
\end{equation*}
By real linearity, the event law is consequently forced to be the
inverse-carried Hermitian congruence. For a normalized Lorentz element
\(\pv{T}\cherm{\pv{T}}=1\), it is
\begin{align}
\pv{X}'
&=\pv{T}^{-1}\pv{X}(\pv{T}^{-1})^{\mathsf H}\notag\\
&=\pv{T}^{*\mathsf H}\pv{X}\pv{T}^*.
\label{eq:general-event-congruence}
\end{align}
Distinct roles are assigned to the involutions. In the Pauli realization,
\({}^{\mathsf H}\) is Hermitian adjunction, whereas
\({}^{*\mathsf H}\) is the adjugate. The latter equals \(\pv{T}^{-1}\) only
after the determinant-one normalization \(\pv{T}\cherm{\pv{T}}=1\) has been imposed.
The real subspace is preserved because
\((\pv{X}')^{\mathsf H}=\pv{X}'\), and the event determinant is preserved because
\(\pv{T}^{-1}\) has unit determinant. When \(\pv{X}\) is invertible, the
transformed inverse is obtained without a new definition:
\begin{equation}
(\pv{X}')^{-1}=\pv{T}^{\mathsf H}\pv{X}^{-1}\pv{T},
\label{eq:event-inverse-transform}
\end{equation}
because the identity is obtained by direct multiplication with
\(\pv{X}'=\pv{T}^{-1}\pv{X}(\pv{T}^{-1})^{\mathsf H}\). The algebraic
pattern for the dual derivative paravector is also supplied by this
complementary law. The event intertwiner is obtained by multiplication of
Eq.\ \eqref{eq:general-event-congruence} on the left by \(\pv{T}\):
\begin{equation}
\pv{T}\pv{X}'=\pv{X}\pv{T}^*.
\label{eq:event-intertwining}
\end{equation}
The form \(\pv{T}\pv{X}'=\pv{X}\pv{T}\) is obtained for a circular rotor because
\(\pv{R}^*=\pv{R}\), but not for a boost. The similarity and event
intertwiners are therefore identical for circular rotors but remain distinct
for boosts.

In contrast, the action of a bivector such as the electromagnetic field is
represented by left multiplication on a carrier. The following similarity is
then required by the original intertwining argument:
\begin{equation}
\pv{F}'=\pv{T}^{-1}\pv{F}\pv{T}.
\label{eq:general-field-similarity}
\end{equation}
For a circular rotor, \(\pv{R}^*=\pv{R}\) and
\(\pv{R}^{-1}=\pv{R}^{\mathsf H}\), so Eq.\ \eqref{eq:general-event-congruence}
is reduced to the active spatial rotation
\begin{equation}
\pv{X}'=\herm{\pv{R}}\pv{X}\pv{R}.
\label{eq:rotor}
\end{equation}
For the hyperbolic element in Eq.\ \eqref{eq:boost-element},
\(\pv{L}^{\mathsf H}=\pv{L}\) and
\(\pv{L}^{-1}=\pv{L}^*=\pv{L}^{*\mathsf H}\), so the same general event law
is reduced to the passive boost
\begin{equation}
\pv{X}'=\conj{\pv{L}}\pv{X}\conj{\pv{L}}.
\label{eq:boost}
\end{equation}
By contrast,
Eq.\ \eqref{eq:general-field-similarity} is reduced to
\(\pv{F}'=\pv{L}^*\pv{F}\pv{L}\).
The unequal right factors are therefore consequences of the representations
being transported, not of an inconsistent choice of inverse.
Under the definition of \(\pv{L}(\theta,\vct{u})\), the sandwich
\(\pv{L}^{\mathsf H}\pv{X}\pv{L}=\pv{L}\pv{X}\pv{L}\) is the inverse-direction event boost, obtained by
replacing \(\theta\) with \(-\theta\). It is physically legitimate, but it
is not the passive \(+\theta\) convention used in Eq.\ \eqref{eq:boost}.

The hyperbolic reduction of the general complex formula is obtained
directly with
\[
f=\cosh(\theta/2),
\qquad
\vct{g}=\sinh(\theta/2)\vct{u},
\qquad
\vct{u}^2=1,
\]
Eq.\ \eqref{eq:similarity-components} is reduced to
\begin{equation}
\vct{V}'
=(\vct{u}\cdot\vct{V})\vct{u}
+\cosh(\theta)\bigl(\vct{V}-(\vct{u}\cdot\vct{V})\vct{u}\bigr)
-i\sinh(\theta)\,\vct{u}\times\vct{V}.
\label{eq:hyperbolic-similarity}
\end{equation}
For \(\vct{V}=\vct{E}+i\vct{B}\), its real and imaginary parts are
exactly the electric--magnetic boost formulas. For a real \(\vct{V}\),
however, the cross term is generally imaginary; hyperbolic similarity is
therefore a bivector or basis transformation rather than an event boost.

The distinction is exposed by the infinitesimal forms. If
\(\pv{T}=1+\pv{K}+O(\pv{K}^2)\), the order functions defined in
Eqs.\ \eqref{eq:com-function} and \eqref{eq:acom-function} give
\begin{align*}
\pv{Z}'&=\pv{Z}+\com{\pv{Z}}{\pv{K}}+O(\pv{K}^2),\\
\pv{X}'&=\pv{X}-\pv{K}\pv{X}-\pv{X}\pv{K}^{\mathsf H}+O(\pv{K}^2).
\end{align*}
For a circular generator \(\pv{K}^{\mathsf H}=-\pv{K}\), the event law is reduced
to the same commutator by which the similarity rotation is generated. For a
hyperbolic generator \(\pv{K}^{\mathsf H}=\pv{K}\), an anticommutator is obtained instead:
\(\pv{X}'=\pv{X}-\acom{\pv{K}}{\pv{X}}+O(\pv{K}^2)\). Spatial vectors are
rotated by the former, while the scalar time component is mixed with the
vector component parallel to the boost by the latter. Thus similarity
rotations are produced by circular generators, whereas event congruences
are produced by hyperbolic generators.

A transform element \(\pv{T}(\lambda)\) is taken to depend on one real
parameter \(\lambda\). All displayed derivatives are ordinary derivatives
with respect to that parameter. The similarity generator is defined by
\begin{equation}
\pv{\Omega}:=\pv{T}^{-1}\frac{\dd\pv{T}}{\dd\lambda}.
\label{eq:similarity-generator}
\end{equation}
Differentiation of the inverse and of the transported paravector then gives
\begin{align*}
\frac{\dd(\pv{T}^{-1})}{\dd\lambda}
&=-\pv{T}^{-1}\frac{\dd\pv{T}}{\dd\lambda}\pv{T}^{-1},\\
\frac{\dd\pv{Z}'}{\dd\lambda} &=
\pv{T}^{-1}\frac{\dd\pv{Z}}{\dd\lambda}\pv{T}
+\com{\pv{Z}'}{\pv{\Omega}},\notag \\
\frac{\dd^2\pv{Z}'}{\dd\lambda^2} &=
\pv{T}^{-1}\frac{\dd^2\pv{Z}}{\dd\lambda^2}\pv{T}\notag \\
&\quad +2\com{\dfrac{\dd\pv{Z}'}{\dd\lambda}}{\pv{\Omega}},\notag \\
&\quad +\com{\pv{Z}'}{\dfrac{\dd\pv{\Omega}}{\dd\lambda}},\notag \\
&\quad +\com{\pv{\Omega}}{\com{\pv{Z}'}{\pv{\Omega}}}.
\end{align*}
They are not space-time derivatives.

\subsubsection{Transformed Basis}
The parameter-dependent element is written as
\begin{align*}
\pv{T}(\lambda)&=f(\lambda)+\vct{g}(\lambda)\cdot\sigv,
\qquad f\in\C,\quad \vct{g}\in\C^3,\\
\pv{T}^{-1}&=\frac{f-\vct{g}\cdot\sigv}{f^2-\vct{g}^2},
\qquad f^2-\vct{g}^2\neq0.
\end{align*}
The transformed basis function is defined by
\begin{equation}
\pv{\sigma}_n':=\pv{T}^{-1}\pv{\sigma}_n\pv{T}.
\label{eq:transformed-basis-definition}
\end{equation}
Its inverse relation is
\[
\pv{\sigma}_n=\pv{T}\pv{\sigma}_n'\pv{T}^{-1}.
\]
Both real hyperbolic elements and imaginary-vector rotors are included by
allowing complex \(\vct{g}\). The following relations are obtained
immediately by similarity:
\begin{align*}
(\pv{\sigma}_n')^2&=1,\\
\pv{\sigma}_n'\pv{\sigma}_m'
&=\delta_{nm}+\sum_{k=1}^3\epsilon_{nmk}\,i\pv{\sigma}_k',\\
\pv{\sigma}_1'\pv{\sigma}_2'\pv{\sigma}_3'&=i.
\end{align*}
Thus the same Clifford algebra is satisfied by the transformed generators
for every invertible \(\pv{T}\).

The separate actions of the two basic involutions on the transformed
generators are represented by
\begin{align*}
\conj{\pv{\sigma}_n'}
&=\conj{(\pv{T}^{-1})}\conj{\pv{\sigma}_n}\conj{\pv{T}}
=-(\conj{\pv{T}})^{-1}\pv{\sigma}_n\conj{\pv{T}},\\
\herm{\pv{\sigma}_n'}
&=\herm{\pv{T}}\herm{\pv{\sigma}_n}\herm{(\pv{T}^{-1})}
=\herm{\pv{T}}\pv{\sigma}_n(\herm{\pv{T}})^{-1}.
\end{align*}
The scalar part is also preserved by similarity, so
\(\eqtext{scalar}(\pv{\sigma}_n')=\eqtext{scalar}(\pv{\sigma}_n)=0\).
The individual star and \(H\) actions are recorded separately by these
formulas and are combined by the adjugate relation.

For every invertible \(\pv{T}\), commutation of the adjugate with the
similarity is obtained from
\(\cherm{\pv{T}}=(f^2-\vct{g}^2)\pv{T}^{-1}\) and
\(\cherm{(\pv{T}^{-1})}=\pv{T}/(f^2-\vct{g}^2)\):
\begin{align*}
(\pv{\sigma}_n')^{*\mathsf H}
&=(f^2-\vct{g}^2)\pv{T}^{-1}(-\pv{\sigma}_n)
  \frac{\pv{T}}{f^2-\vct{g}^2}\\
&=-\pv{T}^{-1}\pv{\sigma}_n\pv{T}=-\pv{\sigma}_n'.
\end{align*}
Consequently a paravector written in the transformed basis,
\[
\pv{Z}=Z_s+\vct{Z}_{\vct{v}}\cdot\sigv',
\]
is assigned the same determinant quadratic form,
\[
\pv{Z}\pv{Z}^{*\mathsf H}=Z_s^2-\vct{Z}_{\vct{v}}^2.
\]

For differentiation with respect to \(\lambda\), the generator from
Eq.\ \eqref{eq:similarity-generator} gives
\begin{align*}
\frac{\dd\pv{\sigma}_n'}{\dd\lambda}
&=\com{\pv{\sigma}_n'}{\pv{\Omega}},&
\frac{\dd\sigv'}{\dd\lambda}
&=\com{\sigv'}{\pv{\Omega}}.
\end{align*}
The coefficient form is obtained directly from
\begin{align*}
\pv{\Omega}
&=\frac{f\,\dfrac{\dd f}{\dd\lambda}
        -\vct{g}\cdot\dfrac{\dd\vct{g}}{\dd\lambda}}
        {f^2-\vct{g}^2}\\
&\quad+
\frac{f\,\dfrac{\dd\vct{g}}{\dd\lambda}
      -\dfrac{\dd f}{\dd\lambda}\vct{g}
      -i\,\vct{g}\times\dfrac{\dd\vct{g}}{\dd\lambda}}
     {f^2-\vct{g}^2}\cdot\sigv.
\end{align*}
Relative to the primed basis, its vector coefficient is written uniquely as
\[
\pv{\Omega}=\Omega_s+\tfrac12(\vct{k}+i\vct{\boldsymbol\omega})\cdot\sigv',
\qquad \vct{k},\vct{\boldsymbol\omega}\in\R^3.
\]
The same generator carried into the unprimed basis is defined by
\begin{equation}
\pv{\Omega}':=\pv{T}\pv{\Omega}\pv{T}^{-1}.
\label{eq:carried-similarity-generator}
\end{equation}
The definition in Eq.\ \eqref{eq:similarity-generator} then gives
\[
\pv{\Omega}'=\frac{\dd\pv{T}}{\dd\lambda}\pv{T}^{-1}.
\]
The coefficients are recovered through scalar extraction.
For \(n\in\{1,2,3\}\), the coefficients are given by cyclicity of the
algebraic scalar part:
\[
k_n+i\omega_n
=2\,\eqtext{scalar}(\pv{\Omega}\pv{\sigma}_n')
=2\,\eqtext{scalar}(\pv{\Omega}'\pv{\sigma}_n).
\]
The scalar part is central with respect to the basis, while
\begin{align*}
\frac{\dd\sigv'}{\dd\lambda}
&=(-\vct{\boldsymbol\omega}+i\vct{k})\times\sigv',\\
\frac{\dd\pv{Z}}{\dd\lambda}
&=\frac{\dd Z_s}{\dd\lambda}+
\left(\frac{\dd\vct{Z}_{\vct{v}}}{\dd\lambda}
      -\vct{Z}_{\vct{v}}\times\vct{\boldsymbol\omega}
      +i\,\vct{Z}_{\vct{v}}\times\vct{k}\right)\cdot\sigv'.
\end{align*}
These are algebraic moving-basis identities. No space-time, energy, or
proper-time interpretation is implied by similarity alone; a specified
Lorentz congruence and frame geometry are required for such an interpretation.

\subsubsection{Rotation}
Equation \eqref{eq:rotor} is obtained as the rotation case of the general
transform law by specializing the transform element to \(\pv{T}=\pv{R}\).
Sigma number two is assigned to the generator, so the exponential is circular
rather than hyperbolic. The rotor function defined in
Eq.\ \eqref{eq:rotor-element} has the expanded form
\begin{align*}
\pv{R}(\theta ,\vct{u}) &= \exp( \tfrac{1}{2} \theta (\vct{u} \cdot i \sigv ) ),\notag \\
&= \cos( \tfrac{1}{2} \theta ) + \sin( \tfrac{1}{2} \theta ) (\vct{u} \cdot i \sigv )\qquad \text{rotor}.
\end{align*}
\begin{equation*}
\pv{R}^{-1} = \exp( -\tfrac{1}{2} \theta (\vct{u} \cdot i \sigv ) ).
\end{equation*}
\(\theta\in\R\), \(\vct{u}\in\R^3\) is a unit real vector, and \(\pv{R}\)
is generated by the sigma-number-two element \(\vct{u}\cdot i\sigv\). The
inverse relation is
\(\pv{R}^{-1}=\pv{R}^{\mathsf H}=\pv{R}^{*\mathsf H}\).
The orientation sign is fixed by the substitution
\begin{equation*}
f=\cos(\theta/2),
\qquad
\vct{g}=i\sin(\theta/2)\vct{u}
\end{equation*}
into Eq.\ \eqref{eq:similarity-components}. Then
\begin{align*}
f^2-\vct{g}^2&=1,\\
f^2+\vct{g}^2&=\cos(\theta),\\
-2(\vct{g}\cdot\vct{r})\vct{g}
&=(1-\cos(\theta))(\vct{u}\cdot\vct{r})\vct{u},\\
-2if(\vct{g}\times\vct{r})
&=\sin(\theta)(\vct{u}\times\vct{r}).
\end{align*}
Therefore the transformed paravector is
\begin{align*}
\pv{X}' &= \pv{R}^{\mathsf H} \pv{X} \pv{R},\notag \\
&= t' + \vct{r}' \cdot \sigv.
\end{align*}
\begin{align*}
t' &= t,\notag \\
\vct{r}' &= \vct{r} \cos(\theta ) + (\vct{u} \times \vct{r}) \sin(\theta ),\notag \\
&\quad + (\vct{u} \cdot \vct{r}) (1 - \cos(\theta )) \vct{u},\notag \\
\qquad \text{Rodrigues's formula}.
\end{align*}
For a parameter-dependent rotor, the specialization of the general generator
is \(\pv{R}^{\mathsf H}(\dd\pv{R}/\dd\lambda)\). Its form carried to the
other side is
\((\dd\pv{R}/\dd\lambda)\pv{R}^{\mathsf H}\).
The following identities are obtained by differentiating
\(\pv{R}^{\mathsf H}\pv{R}=1\):
\[
\left(\pv{R}^{\mathsf H}\frac{\dd\pv{R}}{\dd\lambda}\right)^{\mathsf H}
=-\pv{R}^{\mathsf H}\frac{\dd\pv{R}}{\dd\lambda},
\qquad
\left(\frac{\dd\pv{R}}{\dd\lambda}\pv{R}^{\mathsf H}\right)^{\mathsf H}
=-\frac{\dd\pv{R}}{\dd\lambda}\pv{R}^{\mathsf H}.
\]
The rotation transport law is the specialization of the general
one-parameter similarity formula:
\begin{align*}
\frac{\dd\pv{X}'}{\dd\lambda}
&=\pv{R}^{\mathsf H}\frac{\dd\pv{X}}{\dd\lambda}\pv{R}
+\com{\pv{X}'}{\pv{R}^{\mathsf H}\dfrac{\dd\pv{R}}{\dd\lambda}}\\
&=\pv{R}^{\mathsf H}
\left(\frac{\dd\pv{X}}{\dd\lambda}
+\com{\pv{X}}{\dfrac{\dd\pv{R}}{\dd\lambda}\pv{R}^{\mathsf H}}\right)\pv{R}.
\end{align*}

\subsubsection{Spin Rotors}
Each oriented complementary bivector is identified by the pseudoscalar
without an index-order ambiguity:
\[
i\pv{\sigma}_1=\pv{\sigma}_2\pv{\sigma}_3,\qquad
i\pv{\sigma}_2=\pv{\sigma}_3\pv{\sigma}_1,\qquad
i\pv{\sigma}_3=\pv{\sigma}_1\pv{\sigma}_2.
\]
For unit \(\vct{n},\vct{u}\in\R^3\), the rotor function defined in
Eq.\ \eqref{eq:rotor-element} is evaluated as
\begin{align*}
\pv{R}(\phi,\vct{n})
&=\exp\!\left(\frac{\phi}{2}\,\vct{n}\cdot i\sigv\right)\\
&=\cos(\phi/2)
 +\sin(\phi/2)\,\vct{n}\cdot i\sigv .
\end{align*}
Then
\[
\pv{R}^{-1}=\pv{R}^{\mathsf H}=\pv{R}^{*\mathsf H},\qquad
\pv{R}^{\mathsf H}\pv{R}=1,
\]
and
\[
\pv{R}^{\mathsf H}(\vct{u}\cdot\sigv)\pv{R}
=\vct{u}'\cdot\sigv
\]
with \(\vct{u}'\) defined by Rodrigues's formula. The rotor itself is
single-valued and \(4\pi\)-periodic:
\[
\pv{R}(\phi+2\pi,\vct{n})=-\pv{R}(\phi,\vct{n}),\qquad
\pv{R}(\phi+4\pi,\vct{n})=\pv{R}(\phi,\vct{n}).
\]
The induced spatial rotation is two-to-one because the same conjugation is
produced by \(\pv{R}\) and \(-\pv{R}\).

In the Pauli matrix realization, the two-by-two representation is transposed
and complex-conjugated by \(H\). A general
\(\pv{\psi}\in\C+\C^3\) is therefore a full two-by-two matrix, not one
selected two-entry column: its two columns are rotated independently by
left multiplication. One column is selected by a fixed projector on the
right when a single Pauli wave function is intended.
Under the active convention associated with this vector action,
\[
\pv{\psi}'=\pv{R}^{\mathsf H}\pv{\psi},\qquad
\pv{\psi}'\pv{\psi}'^{\mathsf H}
=\pv{R}^{\mathsf H}(\pv{\psi}\pv{\psi}^{\mathsf H})\pv{R}.
\]
A scalar pairing, normalization, expectation map, projected carrier, and
specified operator action are additionally required for a physical spin state. The
algebra action is thereby separated from the choice of a single physical
spin state.
